\documentclass[11pt,a4paper]{article}

\usepackage[T1]{fontenc}
\usepackage[utf8]{inputenc}
\usepackage{lmodern}
\usepackage{geometry}
\usepackage{setspace}
\usepackage{hyperref}

\geometry{margin=1in}
\onehalfspacing

\title{Off-Chip Memory Support for MASE Hardware Emission\\
\large Coursework Report}
\author{}
\date{\today}

\begin{document}

\maketitle

\section{Introduction}

The MASE hardware emission flow already provides a strong foundation for rapid
translation of quantized neural networks into synthesizable RTL, with automatic
generation of top-level modules, dependencies, and cocotb-based verification
infrastructure. This makes it highly effective for educational and research
prototyping, particularly when targeting compact models that fit comfortably in
on-chip resources.

However, the baseline FPGA-oriented flow is primarily optimized for storing
layer parameters on chip, typically through BRAM-backed sources. For larger
models and low-cost FPGA devices, this becomes a critical bottleneck: available
BRAM limits model size before arithmetic resources are fully utilized. In
practical deployments, especially where price and power constraints are strict,
off-chip parameter storage is often required to scale beyond this limit.

This report presents an extension to the MASE emission and verification flow to
support off-chip compatible designs, with a focus on streaming weights and
biases from external memory into compute modules. The goal is to preserve the
existing structure of the toolchain while adding minimal, targeted changes to
metadata handling, top-level interface generation, and testbench behavior.

Alongside the hardware-generation changes, we introduce a verification path that
drives valid transactions into both activation inputs and externally streamed
parameter ports. This ensures that off-chip behavior is not only representable
in generated RTL, but also testable in a reproducible simulation workflow.

The implementation is evaluated on a use case motivated by collaboration with
UCL's Optical Networks Group, where larger learned classifiers are desirable
but constrained by device memory budgets. This context provides a concrete
motivation for moving from BRAM-only parameter delivery to an off-chip memory
model, and frames the trade-off between achievable model scale and throughput.

\end{document}

